\section{Въведение в статистиката и концепцията за извадки}

Започваме нашето занимание със статистиката — наука, която лежи в основата на вероятностите, използвайки техния фундамент за изграждането на апарат, техники и модели, чрез които можем да изучаваме явления около нас, за които имаме непълна или липсваща информация.

Какво означава това? Представете си, че имате някаква генерална съвкупност — например всички хора в дадена държава, всички клетки в организма, молекули или гласоподаватели. Искаме да изучим някакво тяхно свойство, като например средната височина на хората в България или пропорцията на мъжете и жените. Оказва се обаче, че изследването на всеки един обект от тази съвкупност е твърде скъпо и непрактично. Вместо да разпитваме всеки отделен човек (което би било непосилно), ние правим оценка за свойството на базата на \emph{извадка} — подмножество от обектите, което избираме, за да разберем какво е разпределението на изследваната характеристика.

Самият избор на извадка не е тривиална задача. В историята има показателни примери за провали при съставянето ѝ. Един от тях е през 1936 г. за президентските избори в САЩ. Една компания прави опит да прогнозира изхода от изборите въз основа на много голяма извадка, набирайки я чрез телефонни указатели. По онова време обаче телефони са имали предимно по-богатите хора. Поради това извадката е имала систематично изкривяване (bias) към определена социална страта и прогнозата им се проваля. За разлика от тях, фирмата на Джордж Галъп с едва 4000 респонденти успява да прогнозира правилно, защото тяхната извадка е била представителна и подбрана напълно случайно, отразявайки реалната структура на обществото (включително както по-богати, така и по-бедни). В нашите разглеждания ние винаги ще предполагаме, че извадката се прави напълно случайно и е представителна за генералната съвкупност.

\subsection{Математическа постановка}

Имаме някаква случайна величина $X$, описваща интересуващото ни свойство (например височина, доход, глас за определен кандидат). Тя има функция на разпределение $F_X(x) = \mathbb{P}(X < x)$, а в случай на непрекъсната величина — и плътност $f_X(x)$.

Целта ни е, чрез извадка от $n$ на брой наблюдения, да оценим тази функция на разпределение или нейната плътност. Тъй като изборът е случаен, самите наблюдения $X_1, X_2, \dots, X_n$ формират вектор от \emph{независими и еднакво разпределени случайни величини} (н.е.р.с.в.), като всяко $X_j$ има същото разпределение като $X$:
\[
    \vec{X} = (X_1, X_2, \dots, X_n)
\]
Това е концептуалната постановка. В практиката ние провеждаме експеримента или запитването и получаваме конкретни числови стойности $x_1, x_2, \dots, x_n$, които представляват реализацията на случайната величина $\vec{X}$. От тези числа, ние трябва да направим обосновано предположение за функцията на разпределение $F_X(x)$ или за плътността $f_X(x)$.

\section{Точкови оценки}

Ако нямаме никакви предварителни допускания за функцията на разпределение на $X$, задачата е много трудна. В класическата параметрична статистика ние стесняваме задачата, като допускаме, че $X$ принадлежи на определен \emph{клас} от случайни величини, параметризиран с даден параметър $\theta$.

Например:
\begin{itemize}
    \item $X \sim \Ber(p)$ — Бернулиево разпределение, където $\theta = p$ (вероятността за успех, напр. гласуване за даден кандидат).
    \item $X \sim \text{Exp}(\lambda)$ — Експоненциално разпределение за време на живот на електронни компоненти, където $\theta = \lambda$.
    \item $X \sim N(\mu, \sigma^2)$ — Нормално разпределение, където параметърът е двумерен вектор $\theta = (\mu, \sigma^2)$.
\end{itemize}

По този начин, вместо да търсим произволна функция, нашата цел се свежда до това да намерим оценка $\hat{\theta}$ за параметъра $\theta$ въз основа на наблюденията. Това води до концепцията за \emph{точкова оценка} (Point Estimator).

Точковата оценка $\hat{\theta} = \hat{\theta}(\vec{X})$ е функция на случайния вектор от наблюденията, която ни служи като приближение за истинската, но неизвестна стойност на параметъра $\theta$. Тъй като всяка функция може формално да се разглежда като оценка, имаме нужда от методи за конструиране на оценки, които притежават оптимални свойства. Два от най-популярните метода са Методът на максималното правдоподобие (МПО) и Методът на моментите (ММО).

\section{Метод на максималното правдоподобие (МПО)}

Методът на максималното правдоподобие, или Maximum Likelihood Estimation (MLE), почива на една много интуитивна идея: от всички възможни стойности за параметъра $\theta$ ние трябва да изберем тази, при която вероятността (или съвместната плътност) да наблюдаваме данните, които реално сме получили, е възможно най-голяма.

\subsection{Дефиниция на функция на правдоподобие}
Нека $X$ е случайна величина, чиято функция на разпределение (или плътност) $f_X(x; \theta)$ е от клас, параметризиран с $\theta$. Имаме независимата извадка $\vec{X} = (X_1, X_2, \dots, X_n)$.

\begin{keydefn}[функция на правдоподобие]\label{def:likelihood}
\emph{Функция на правдоподобие} $L(\vec{X}, \theta)$ наричаме съвместната плътност (или вероятност в дискретния случай) на вектора $\vec{X}$. Тъй като наблюденията са независими, съвместната плътност се разлага в произведение от индивидуалните плътности:
\[
    L(\vec{X}, \theta) = \prod_{j=1}^n f_X(X_j; \theta) .
\]
\end{keydefn}

За конкретна реализация на извадката $\vec{x} = (x_1, \dots, x_n)$, функцията на правдоподобие от дефиниция~\ref{def:likelihood} е:
\[
    L(\vec{x}, \theta) = \prod_{j=1}^n f_X(x_j; \theta)
\]

\textbf{Максимално правдоподобна оценка} (МПО) $\hat{\theta}$ е тази стойност на $\theta$, която максимизира функцията на правдоподобие:
\[
    L(\vec{X}, \hat{\theta}) = \sup_{\theta} L(\vec{X}, \theta)
\]

За удобство при пресмятанията, понеже логаритъмът е строго растяща функция, максимизирането на $L(\vec{X}, \theta)$ е еквивалентно на максимизирането на логаритъма на функцията на правдоподобие:
\[
    \ln L(\vec{X}, \theta) = \sum_{j=1}^n \ln f_X(X_j; \theta)
\]
Ако $f_X(x; \theta)$ е диференцируема по $\theta$, можем да намерим екстремума, решавайки системата уравнения:
\[
    \frac{\partial \ln L(\vec{X}, \theta)}{\partial \theta} = 0
\]

\begin{example}[Равномерно разпределение $U(0, \theta)$]\label{ex:12-1}

Допускаме, че $X$ е равномерно разпределена в интервала $(0, \theta)$, т.е. $X \sim U(0, \theta)$. Параметърът $\theta > 0$. Плътността е дадена чрез индикаторната функция:
\[
    f_X(x; \theta) = \frac{1}{\theta} \ind_{(0, \theta)}(x)
\]
Функцията на правдоподобие за $n$ наблюдения е:
\[
    L(\vec{X}, \theta) = \prod_{j=1}^{n} \frac{1}{\theta} \ind_{(0, \theta)}(X_j) = \frac{1}{\theta^n} \ind_{(0, \theta)}(X^*) = \begin{cases} \frac{1}{\theta^n} & \text{ако } X^* \le \theta \\ 0 & \text{ако } X^* > \theta \end{cases}
\]
където $X^* = \max_{j \le n}(X_j)$ е максималната стойност сред наблюденията. Забележете, че произведението на индикаторните функции $\prod_{j=1}^n \ind_{(0, \theta)}(X_j)$ е равно на 1 тогава и само тогава, когато всички $X_j < \theta$, което е еквивалентно на това техният максимум $X^*$ да е по-малък от $\theta$.

Търсим $\hat{\theta}$, което максимизира $L(\vec{X}, \theta)$:
\[
    L(\vec{X}, \hat{\theta}) = \sup_{\theta > 0} L(\vec{X}, \theta)
\]
Тъй като функцията $\frac{1}{\theta^n}$ е строго намаляваща по $\theta$, нейната максимална стойност се достига в най-малката възможна точка за $\theta$, при която функцията не е нула. Според индикатора, трябва $\theta \ge X^*$. Следователно най-голямата стойност се постига, когато:
\[
    \hat{\theta}_1 = X^* = \max_{j \le n}(X_j)
\]
Ако например сме получили наблюденията 1, 2, 3, 4 и 5, максимално правдоподобната оценка за $\theta$ би била просто най-голямата стойност — 5.

\end{example}

\begin{example}[Нормално разпределение $N(\mu, \sigma^2)$]\label{ex:12-2}

Нека $X \sim N(\mu, \sigma^2)$. Имаме два параметъра за оценяване: $\theta = (\mu, \sigma^2)$.
Плътността на нормалното разпределение за едно наблюдение е:
\[
    f_X(x; \mu, \sigma^2) = \frac{1}{\sqrt{2\pi\sigma^2}} e^{-\frac{(x - \mu)^2}{2\sigma^2}}
\]
Съставяме функцията на правдоподобие за извадката $\vec{X}$:
\[
    L(\vec{X}; \theta) = \left(\frac{1}{\sqrt{2\pi}}\right)^n \sigma^{-n} e^{-\sum_{j=1}^n \frac{(X_j - \mu)^2}{2\sigma^2}}
\]
Логаритмуваме я, за да превърнем произведението в сума, което значително улеснява диференцирането:
\[
    \ln L(\vec{X}, \theta) = n \ln \frac{1}{\sqrt{2\pi}} - \frac{n}{2} \ln \sigma^2 - \sum_{j=1}^n \frac{(X_j - \mu)^2}{2\sigma^2}
\]
Намираме частната производна спрямо $\mu$ и я приравняваме на нула:
\[
    \frac{\partial \ln L}{\partial \mu} = \frac{1}{\sigma^2} \sum_{j=1}^n (X_j - \mu) = 0
\]
\[
    \sum_{j=1}^n X_j - n\mu = 0 \implies \hat{\mu} = \frac{1}{n} \sum_{j=1}^n X_j = \overline{X}_n
\]
Оценката за $\mu$ е средното аритметично на наблюденията. Забележете, че тя не зависи от това дали познаваме $\sigma^2$ или не.

Сега намираме частната производна спрямо $\sigma^2$ и я приравняваме на нула:
\[
    \frac{\partial \ln L}{\partial \sigma^2} = -\frac{n}{2\sigma^2} + \frac{1}{2\sigma^4} \sum_{j=1}^n (X_j - \mu)^2 = 0
\]
Ако параметърът $\mu$ ни е предварително известен, решението директно ни дава:
\[
    \hat{\sigma}^2 = \frac{1}{n} \sum_{j=1}^n (X_j - \mu)^2
\]
Ако $\mu$ е неизвестен, заместваме неговата оценка $\hat{\mu} = \overline{X}_n$ в уравнението, откъдето получаваме:
\[
    \hat{\sigma}^2 = \frac{1}{n} \sum_{j=1}^n (X_j - \overline{X}_n)^2
\]


\end{example}

\section{Метод на моментите (ММО)}

Методът на моментите почива на Закона за големите числа, според който емпиричните (извадкови) моменти се схождат почти сигурно към теоретичните моменти на разпределението при увеличаване на обема на извадката.

Нека въведем емпиричния момент от ред $j$:
\[
    \bar{X}_n^{(j)} = \frac{1}{n} \sum_{k=1}^n X_k^j, \quad j \ge 1
\]
Това е средното аритметично от наблюденията, повдигнати на $j$-та степен. (При $j=1$ получаваме обикновеното средно аритметично $\bar{X}_n$).

\begin{defn}[ММО]
Нека $X$ е случайна величина с функция на разпределение $F_X(x; \theta)$, където параметърът е $s$-мерен вектор $\theta = (\theta_1, \dots, \theta_s)$. Оценката по метода на моментите $\hat{\theta} = (\hat{\theta}_1, \dots, \hat{\theta}_s)$ намираме, решавайки следната система от $s$ уравнения:
\[
    \mu^{(j)}(\theta) = \bar{X}_n^{(j)}, \quad 1 \le j \le s
\]
където $\mu^{(j)}(\theta) = \E[X^j]$ са теоретичните моменти на случайната величина като функция на параметъра $\theta$.
\end{defn}

Идеята е да приравним първите $s$ на брой теоретични момента към съответните извадкови моменти и да решим получената система спрямо неизвестните параметри. Тъй като $\bar{X}_n^{(j)} \to \E[X^j]$ при $n \to \infty$, методът дава оценки, които при определени условия на обратимост клонят към истинските стойности на параметрите.

\begin{example}[Равномерно разпределение $U(0, \theta)$ (ММО)]\label{ex:12-3}

Разглеждаме отново случая $X \sim U(0, \theta)$. Тъй като имаме само един параметър, се нуждаем от едно уравнение — за първия момент.
Теоретичното очакване на $X$ е:
\[
    \E[X] = \frac{\theta}{2}
\]
Емпиричният първи момент е средното аритметично $\bar{X}_n$. Приравняваме ги:
\[
    \frac{\theta}{2} = \bar{X}_n
\]
Оттук директно изразяваме оценката по ММО за $\theta$:
\[
    \hat{\theta}_2 = 2 \bar{X}_n
\]
Виждаме, че за равномерното разпределение двата метода (МПО и ММО) дават коренно различни оценки ($\hat{\theta}_1 = \max(X_j)$ срещу $\hat{\theta}_2 = 2\bar{X}_n$).

\end{example}

\begin{example}[Нормално разпределение $N(\mu, \sigma^2)$ (ММО)]\label{ex:12-4}

Нека $X \sim N(\mu, \sigma^2)$. Имаме два параметъра ($\mu$ и $\sigma^2$), затова ни трябват първите два момента.
Теоретичните моменти са:
\[
    \E[X] = \mu
\]
\[
    \E[X^2] = \Var(X) + (\E[X])^2 = \sigma^2 + \mu^2
\]
Съставяме системата от уравнения, приравнявайки ги към емпиричните моменти:
\begin{align*}
    \mu &= \bar{X}_n^{(1)} \\
    \sigma^2 + \mu^2 &= \bar{X}_n^{(2)}
\end{align*}
От първото уравнение веднага получаваме:
\[
    \hat{\mu} = \bar{X}_n^{(1)} = \frac{1}{n} \sum_{j=1}^n X_j
\]
Замествайки $\hat{\mu}$ във второто уравнение, получаваме оценката за дисперсията:
\[
    \hat{\sigma}^2 = \bar{X}_n^{(2)} - \left(\bar{X}_n^{(1)}\right)^2 = \frac{1}{n} \sum_{j=1}^n X_j^2 - \left(\frac{1}{n} \sum_{j=1}^n X_j\right)^2 = \frac{1}{n} \sum_{j=1}^n \left(X_j - \bar{X}_n^{(1)}\right)^2
\]
В случая на нормалното разпределение, оценките по метода на моментите (ММО) и по метода на максималното правдоподобие (МПО) съвпадат напълно.

\end{example}

\section{Свойства на точковите оценки}

За да определим коя от получените оценки е по-добра, изследваме техните статистически свойства, като \emph{неизместеност} (unbiasedness) и \emph{състоятелност} (consistency).

\begin{defn}[неизместена оценка]\label{def:unbiased}
Една оценка $\hat{\theta}$ се нарича \emph{неизместена}, ако нейното математическо очакване съвпада с истинската стойност на параметъра:
\[ \E[\hat{\theta}] = \theta . \]
\end{defn}

\begin{supp}[оценка с минимална дисперсия; ефективна оценка]\label{supp:efficient}
Неизместената оценка $\hat{\theta}$ се нарича \emph{оценка с минимална дисперсия},
ако за всяка друга неизместена оценка $\tilde{\theta}$ е изпълнено
$\Var \hat{\theta} \le \Var \tilde{\theta}$. Неизместена оценка с минимална
дисперсия се нарича \emph{ефективна}. Ако за даден параметър съществува ефективна
оценка, тя е единствена.
\end{supp}

\subsection{Свойства на оценките при $U(0, \theta)$}

Сравняваме двете получени оценки: $\hat{\theta}_1 = X^*$ (от МПО) и $\hat{\theta}_2 = 2 \bar{X}_n$ (от ММО).

За оценката по метода на моментите $\hat{\theta}_2$:
\[
    \E[\hat{\theta}_2] = \E[2 \bar{X}_n] = \frac{2}{n} \sum_{j=1}^n \E[X_j] = \frac{2}{n} \cdot n \cdot \frac{\theta}{2} = \theta
\]
Тъй като $\E[\hat{\theta}_2] = \theta$, оценката $\hat{\theta}_2$ е \textbf{неизместена} по смисъла на дефиниция~\ref{def:unbiased}. Освен това, благодарение на Закона за големите числа, $\bar{X}_n \xrightarrow{\text{п.с.}} \frac{\theta}{2}$, следователно $\hat{\theta}_2$ е и \textbf{силно състоятелна} оценка.

За оценката по метода на максималното правдоподобие $\hat{\theta}_1 = X^*$:
За да намерим очакването на $X^*$, първо намираме нейната функция на разпределение $F_{X^*}(x)$ за $x \in (0, \theta)$:
\[
    F_{X^*}(x) = \mathbb{P}(X^* < x) = \mathbb{P}\left(\bigcap_{j=1}^n \{X_j < x\}\right)
\]
Тъй като наблюденията са независими и еднакво разпределени, вероятността от сечението се разпада в произведение от вероятностите:
\[
    F_{X^*}(x) \stackrel{\text{незав.}}{=} \prod_{j=1}^n \mathbb{P}(X_j < x) = \left( \mathbb{P}(X_1 < x) \right)^n = \left(\frac{x}{\theta}\right)^n, \quad x \in (0, \theta)
\]
Плътността се получава чрез диференциране:
\[
    f_{X^*}(x) = n \frac{x^{n-1}}{\theta^n}, \quad x \in (0, \theta)
\]
Сега пресмятаме математическото очакване на $X^*$:
\[
    \E[X^*] = \int_0^\theta x \cdot n \frac{x^{n-1}}{\theta^n} dx = \frac{n}{\theta^n} \left[ \frac{x^{n+1}}{n+1} \right]_0^\theta = \frac{n}{n+1} \theta
\]
Тъй като $\E[\hat{\theta}_1] = \frac{n}{n+1} \theta \neq \theta$, оценката по метода на максималното правдоподобие не изпълнява дефиниция~\ref{def:unbiased} и е \textbf{изместена}. Това е интуитивно логично — максималната стойност в извадката винаги ще бъде малко по-малка от същинската горна граница $\theta$. С увеличаването на извадката обаче, този максимум се доближава все повече до истинското $\theta$ (оценката е асимптотично неизместена). Ако желаем строго неизместена оценка, базирана на $X^*$, можем да я конструираме като:
\[
    \tilde{\theta}_1 = \frac{n+1}{n} X^*
\]
Въпреки че $X^*$ е изместена, тя все пак е \textbf{състоятелна}. Това лесно може да се докаже, като изследваме вероятността $X^*$ да остане отдалечена от $\theta$ на разстояние $\epsilon$:
\[
    \mathbb{P}(X^* < \theta - \epsilon) = \left(\frac{\theta - \epsilon}{\theta}\right)^n
\]
При $n \to \infty$, тъй като дробта в скобите е строго по-малка от 1, този израз клони към 0, което показва, че $X^*$ се схожда по вероятност към $\theta$.

\subsection{Изместеност на дисперсията при Нормално разпределение}

Нека обърнем внимание и на получената оценка за дисперсията при $N(\mu, \sigma^2)$:
\[
    \hat{\sigma}^2 = \frac{1}{n} \sum_{j=1}^n \left(X_j - \bar{X}_n^{(1)}\right)^2
\]
Може да се докаже, че $\E[\hat{\sigma}^2] = \frac{n-1}{n} \sigma^2 \neq \sigma^2$. Тоест, получената по методите МПО и ММО оценка $\hat{\sigma}^2$ е изместена (тя систематично подценява истинската дисперсия, понеже използва емпиричното средно вместо истинското). За да се коригира това изкривяване, в практиката се използва \emph{коригираната извадкова дисперсия}:
\[
    S^2 = \frac{n}{n-1} \hat{\sigma}^2 = \frac{1}{n-1} \sum_{j=1}^{n} (X_j - \bar{X}_n^{(1)})^2
\]
За нея вече е в сила:
\[
    \E[S^2] = \frac{n}{n-1} \E[\hat{\sigma}^2] = \frac{n}{n-1} \frac{n-1}{n} \sigma^2 = \sigma^2
\]
което я прави неизместена оценка за параметъра $\sigma^2$.

\section{Задачи}

\begin{enumerate}
    \item За нормалния модел от пример~\ref{ex:12-2} пресметнете частните производни на $\ln L(\vec{X};\mu,\sigma^2)$ по $\mu$ и по $\sigma^2$, приравнете ги на нула и решете получената система.

    \item Нека $X_1,\dots,X_n$ са независими и равномерно разпределени в $(0,\theta)$ и $X^*=\max_{j\le n}X_j$. Проверете състоятелността на оценката $\hat\theta_1=X^*$, като за $\varepsilon>0$ пресметнете
    \[
    \mathbb{P}(X^*<\theta-\varepsilon)
    \]
    и изследвате границата при $n\to\infty$.
\end{enumerate}
